\documentclass{beamer}
\usepackage[utf8]{inputenc}
\usepackage{amsmath}
\usepackage{graphicx}
\usepackage{hyperref}

\usepackage[T1]{fontenc}
\usepackage[french]{babel}


\usepackage{amssymb}
\usepackage{array}
\usepackage{tikz}
\usepackage{amsmath}
\usepackage{amsmath,amssymb}
\usepackage{graphicx}
\title{Méthode actuarielle }
\subtitle{L'actuariat et l'intelligence Artificielle dans l'UEMOA à l'horizon 2030}
\author{Samuel AYANOU}
\institute{Ecole Nationale de la Statistique et de l'Analyse Economique\\
\textbf{Superviseur :} M. Babacar NDIAYE}

\date{Juin 2025}


\usetheme{Warsaw}
\usecolortheme{whale}
\setbeamertemplate{navigation symbols}{}
\setbeamertemplate{caption}[numbered]
\AtBeginSection[]{
	\begin{frame}
		\frametitle{Plan}
		\tableofcontents[currentsection]
	\end{frame}
}
\begin{document}

\frame{\titlepage}

\begin{frame}
\frametitle{Sommaire}
\tableofcontents
\end{frame}


% INTRODUCTION
\section{Introduction}
\begin{frame}{Introduction Générale}
	\begin{block}{Cadre mondial}
		Transformation numérique des fonctions techniques dans la finance et l’assurance.
	\end{block}
	\begin{block}{Enjeux UEMOA}
		Rattrapage technologique, faiblesse du tissu assurantiel, opportunités démographiques.
	\end{block}
	\begin{block}{Problématique}
		Comment l’IA, et plus spécifiquement l’IA générative, peut-elle transformer le métier d’actuaire dans l’espace UEMOA, tout en contournant les contraintes structurelles ?
	\end{block}
\end{frame}
\section{État des lieux du secteur}
\begin{frame}{État des lieux du secteur}
	\begin{block}{Marché de l’assurance (2022)}
		\begin{itemize}
			\item Primes : 1 223 Mds FCFA , soit ~1,1\% du PIB cumulé des pays membres de l’UEMOA.
			\item Taux de pénétration en moyenne : <2\% (Afrique du Sud : 9\%, France : 12\%)
			\item Assurance vie : 24\% des primes (OCDE : >60\%)
			\item +60\% des contrats = assurance auto obligatoire
		\end{itemize}
	\end{block}
\end{frame}

\begin{frame}{État des lieux du secteur}
	\begin{block}{Démographie et digitalisation}
	\begin{itemize}
		\item 123 millions d’habitants (2023)
		\item 60\% ont moins de 25 ans
		\item >95\% de pénétration mobile (ARCEP Bénin, 2022)
		\item 64\% des adultes ont un compte mobile (BCEAO, 2022)
	\end{itemize}
	\end{block}
	\begin{block}{Capacités actuarielles}
		\begin{itemize}
		\item Peu de formations locales spécialisées
		\item Forte dépendance aux cabinets étrangers
		\item  $<$ à 300 actuaires certifiés dans l’UEMOA (FANAF, 2023)

\end{itemize}
			\end{block}
\end{frame}


\section{Les technologies émergentes}
\begin{frame}{ Les technologies émergentes}
	
	\begin{block}{ IA (LLMs)}
		Large Language Models  (LLMs) capable de générer texte, code, images .\\
		Applications : rédaction, modélisation, automatisation.
	\end{block}
	
	\begin{block}{Vision \& Images}
		Traitement d’images pour télé-expertise et cartographie des risques climatiques.\\
		Évaluation rapide et précise des risques.
	\end{block}
	
	
	
	
\end{frame}

\begin{frame}{L’IA Générative au service de l’actuariat}
	
	\begin{block}{ IoT \& Blockchain}
		IoT : données en temps réel pour affiner la tarification des risques.\\
		Blockchain : contrats automatisés et données fiables pour l’actuaire.\\
		Applications : assurance paramétrique, télématique, smart contracts.
	\end{block}
	\begin{tabular}{|l|l|l|}
		\hline
		Fonction & Application & Gain estimé \\
		\hline
		Tarification & Variables, automatisation & -50\% temps \\
		Modélisation & Hyperparamètres, diagnostic & -60\% erreurs \\
		Reporting & Traduction, mise en forme & -70\% délais \\
		Formation & Tuteurs IA, quizz & +100\% apprentissage \\
		\hline
	\end{tabular}
\end{frame}

\section{Contrainte}
% Contraintes
\begin{frame}{Contraintes structurelles}
	\begin{block}{ Manque de talents}
		1 actuaire / 400 000 hab. (France : 1 / 10 000), fuite des cerveaux.
	\end{block}
	\begin{block}{Faiblesse des données}
		Données incomplètes, silotées, crainte du partage.
	\end{block}
	\begin{block}{ Régulation rigide}
		Code CIMA peu adapté, absence de sandbox, risques éthiques.
	\end{block}
\end{frame}

% Scénarios
\section{Recommandation}
\begin{frame}{Scénarios Prospectifs 2025–2030}
	\begin{block}{Trois trajectoires possibles}
		\begin{itemize}
			\item \textbf{Scénario 1 – Inertie} : Faible investissement dans l’IA, recours accru à des compétences et technologies étrangères, accentuation des écarts avec les marchés matures.
			\item \textbf{Scénario 2 – Transition assistée} : Adoption progressive via des partenariats (startups, bailleurs, ONG), montée en compétence progressive, mais dépendance partielle maintenue.
			\item \textbf{Scénario 3 – Révolution numérique} : Intégration massive et souveraine de l’IA, formation d’actuaires locaux maîtrisant l’IA, et structuration d’un écosystème régional d’innovation.
		\end{itemize}
	\end{block}
\end{frame}

% Feuille de route
\begin{frame}{Feuille de Route Stratégique pour l’UEMOA}
	\begin{block}{Formation et développement des talents}
		\begin{itemize}
			\item Création d’un Institut régional IA–Actuariat
			\item Bourses ciblées et bootcamps pratiques
			\item Partenariats avec universités et entreprises du secteur
		\end{itemize}
	\end{block}
	
	\begin{block}{Infrastructure et gouvernance des données}
		\begin{itemize}
			\item Mise en place d’un Data Commons régional (bases mutualisées)
			\item Génération de données synthétiques pour la modélisation
			\item Protocoles d’anonymisation automatique
		\end{itemize}
	\end{block}
\end{frame}

\begin{frame}{Feuille de Route (suite)}
	\begin{block}{Cadre réglementaire agile}
		\begin{itemize}
			\item Création d’un CIMA Lab IA (espace d’expérimentation réglementaire)
			\item Publication de lignes directrices sur l’équité, la transparence et l’auditabilité algorithmique
		\end{itemize}
	\end{block}
	
	\begin{block}{Soutien à l’innovation locale}
		\begin{itemize}
			\item Incubateurs régionaux pour jeunes actuaires et startups IA
			\item Appels à projets ciblés : micro-assurance, climat, agriculture
			\item Accès au financement via fonds d’innovation panafricains
		\end{itemize}
	\end{block}
\end{frame}


% Conclusion
\section{Conclusion}
\begin{frame}{Conclusion}
	\begin{block}{}
	\begin{itemize}
		\item L’IA = opportunité de rattrapage structurel pour l’UEMOA
		\item L’actuariat : levier d’inclusion financière
		\item Appel à une action collective : public, privé, université, régulation
		\item Vision 2030 : actuaire augmenté, secteur transformé
	\end{itemize}
	\end{block}
\end{frame}


\end{document}